% page 455 of the facsimile (image p-454 of the scan)
% block 162.6 x 224.7 mm ; left margin 39.4 mm
\pgc{17.780mm}{0.000mm}{
\l{\cel{53}447}
\l{}
\l{\sou{plevrito}. (\sou{patol.}) Inflamuro di la plevro. - DEFIRS.}
\l{}
\l{\sou{plexo.(anat.})\cel{2}Reto de fibreti nerva o de vaskuli interplek-}
\l{\cel{3}tita. - DEFIS.}
\l{}
\l{\sou{pleyado.}\cel{2}I. (\sou{astron.}) Singla ek la sis steli (en la epoki an-}
\l{\cel{3}tiqua, onu dicis : sep) qui formacas grupo en la Tauro-ste-}
\l{\cel{3}laro. - II. (\sou{literaturo}) Grupo de sep poeti (en Alexandria}
\l{\cel{3}ed en la Francia di la yarcento 16-esma.) III. (\sou{filoz.})}
\l{\cel{3}Grupo ek la sep saji di la Grekia antiqua. IV. (en\cel{1}\sou{la domeno}}
\l{\cel{3}di\cel{1}\sou{la linguo universala, Ido}).Grupo ek la sep homa fonti}
\l{\cel{3}di la linguo universala definitiva : Louis de Beaufront,}
\l{\cel{3}Louis Couturat, Otto Jespersen, André Lalande, Wilhelm Ost-}
\l{\cel{3}wald, R. Lorenz e L. Pfaundler.- DEFIS.}
\l{}
\l{\sou{plezar}. (\sou{netrans., ad, per, pro}).(\sou{aludante la kozi o la personi})}
\l{\cel{3}Atraktar (\sou{metafore}) per sua charmo. - EFI.}
\l{}
\l{\sou{plezuro}. Sentimento di komforto psikologiala, qua naskas de la}
\l{\cel{3}satisfaceso di deziro, di aspiro. - DEFIS.}
\l{}
\l{\sou{pliko}. (\sou{patol.)}\cel{2}Exuduro ek la pelo haroza, qua aglutinas la}
\l{\cel{3}hari aden fasketi. - EFI.}
\l{}
\l{\sou{plinto.}\cel{2}(\sou{arkitekt}.)I. Tabulo quadrata qua formacas soklo. -}
\l{\cel{3}II. Bendo plata qua jacas an la pedo di edifico, infre di}
\l{\cel{3}muro di apartamento. - III. Bendo de ligno qua aplikesis}
\l{\cel{3}cirkum panelo. - DEFIS.}
\l{}
\l{\sou{plioceno.}- (\sou{geol.}) Strato supra di la sedimento terciara,}
\l{\cel{3}qua kontenas la fosili maxim recenta. - DEFIRS.}
\l{}
\l{\sou{plisar}. (\sou{trans.}) Dispozar tale ke lu formacas plura falduri.-DF}
\l{}
\l{\sou{plombo}. (\sou{kemio}). Metalo bluatre-blanka, mola, tre pezoza. -}
\l{\cel{3}\sou{Simbolo kemiala}\cel{1}: Pb. - DEFIRS.}
\l{}
\l{\sou{pliz.}. Uzata kande onu deziras demandar polite ad ulu, ke}
\l{\cel{3}lu....+ez. - E. (Plu konciza kam \textquotedbl{}Me pregas vu pri lo :...ez\textquotedbl{}}
\l{\cel{3}o \textquotedbl{}Voluntez esar tante komplezema ke vu....us...\textquotedbl{}.)}
\l{}
\l{\sou{plorar(netrans.}) I. Manifestar, expresar sua aflikteso. - II.}
\l{\cel{3}Emisar lakrimi- F.}
\l{}
\l{\sou{plu.}\cel{2}Indikas la augmento o la kresko di quanto o qualeso}
\l{\cel{3}qui duras esar od egardesar kom identa a su ipsa. - FI.}
\l{}
\l{\sou{plugar.}\cel{1}(trans.)\cel{2}Apertar kultivo-sulo, e renversar la glebi,}
\l{\cel{3}per exkavar sulki en tero, uzante la soko di instrumento}
\l{\cel{3}specala. - DR.}
}
